\textbf{Phần 1. Trắc nghiệm lựa chọn}
\setcounterpageref{ex}{0}

\begin{ex} %[2D8V44]
Kết quả của phép chia $35 : 5$ là
\choice
{$6$}
{\True $7$}
{$8$}
{$9$}
\loigiai{
$35 : 5 = 7$.
\textbf{Đáp án của câu hỏi là B}
}
\end{ex}

\begin{ex} %[2D8V44]
Có $40$ quyển vở chia đều cho $5$ tổ học sinh. Mỗi tổ nhận được bao nhiêu quyển vở?
\choice
{$6$ quyển}
{$7$ quyển}
{$9$ quyển}
{\True $8$ quyển}
\loigiai{
$40 : 5 = 8$ (quyển).
\textbf{Đáp án của câu hỏi là D}
}
\end{ex}

\begin{ex} %[2D8V44]
Tính giá trị của biểu thức: $45 : 5 - 4$.
\choice
{\True $5$}
{$4$}
{$6$}
{$7$}
\loigiai{
$45 : 5 = 9$; $9 - 4 = 5$.
\textbf{Đáp án của câu hỏi là A}
}
\end{ex}

\begin{ex} %[2D8V44]
Một đàn gà con có $25$ con được chia đều vào $5$ chuồng. Mỗi chuồng có bao nhiêu con gà?
\choice
{$4$ con}
{$6$ con}
{\True $5$ con}
{$3$ con}
\loigiai{
$25 : 5 = 5$ (con).
\textbf{Đáp án của câu hỏi là C}
}
\end{ex}

\begin{ex} %[2D8V44]
Điền số thích hợp: $50 : 5 = \dots \times 2$.
\choice
{$4$}
{\True $5$}
{$6$}
{$3$}
\loigiai{
Vế trái: $50 : 5 = 10$. Ta có: $10 : 2 = 5$. Vậy điền 5.
\textbf{Đáp án của câu hỏi là B}
}
\end{ex}

\begin{ex} %[2D8V44]
Tìm $x$ biết: $x \times 5 = 30$.
\choice
{$5$}
{$7$}
{$8$}
{\True $6$}
\loigiai{
$x = 30 : 5 = 6$.
\textbf{Đáp án của câu hỏi là D}
}
\end{ex}

\begin{ex} %[2D8V44]
Mẹ có $15$ chiếc bánh, mẹ chia cho các con mỗi người $5$ chiếc bánh. Hỏi mẹ đã chia bánh cho mấy người con?
\choice
{\True $3$ người}
{$2$ người}
{$4$ người}
{$5$ người}
\loigiai{
$15 : 5 = 3$ (người).
\textbf{Đáp án của câu hỏi là A}
}
\end{ex}

\begin{ex} %[2D8V44]
Bác nông dân có $20\text{ kg}$ đỗ xanh đóng đều vào các túi loại $5\text{ kg}$. Hỏi bác đóng được bao nhiêu túi đỗ?
\choice
{$3$ túi}
{$5$ túi}
{\True $4$ túi}
{$6$ túi}
\loigiai{
$20 : 5 = 4$ (túi).
\textbf{Đáp án của câu hỏi là C}
}
\end{ex}

\begin{ex} %[2D8V44]
Điền dấu thích hợp: $30 : 5 \dots 12 : 2$.
\choice
{$>$}
{\True $=$}
{$<$}
{Không so sánh được}
\loigiai{
$30 : 5 = 6$; $12 : 2 = 6$. Hai vế bằng nhau nên điền dấu $=$.
\textbf{Đáp án của câu hỏi là B}
}
\end{ex}

\begin{ex} %[2D8V44]
Tìm một số biết rằng lấy số đó nhân với $5$ rồi trừ đi $15$ thì được $20$.
\choice
{$6$}
{$8$}
{$9$}
{\True $7$}
\loigiai{
Tính ngược: $20 + 15 = 35$; $35 : 5 = 7$.
\textbf{Đáp án của câu hỏi là D}
}
\end{ex}

\begin{ex} %[2D8V44]
Một tấm vải dài $45\text{ m}$ được cắt thành các mảnh dài $5\text{ m}$. Hỏi người thợ cắt được bao nhiêu mảnh vải?
\choice
{\True $9$ mảnh vải}
{$8$ mảnh vải}
{$10$ mảnh vải}
{$7$ mảnh vải}
\loigiai{
$45 : 5 = 9$ (mảnh vải).
\textbf{Đáp án của câu hỏi là A}
}
\end{ex}

\begin{ex} %[2D8V44]
Trong vườn có $50$ cây hoa trồng thành $5$ luống bằng nhau. Mỗi luống có bao nhiêu cây hoa?
\choice
{$8$ cây}
{$9$ cây}
{\True $10$ cây}
{$12$ cây}
\loigiai{
$50 : 5 = 10$ (cây).
\textbf{Đáp án của câu hỏi là C}
}
\end{ex}

\textbf{Phần 2. Trắc nghiệm đúng sai}
\setcounterpageref{ex}{0}

\begin{ex} %[2D8V44]
Xét tính đúng sai của các phép tính trong bảng chia 5:
\choiceTF
{\True Các số bị chia trong bảng chia 5 đều có chữ số tận cùng là 0 hoặc 5.}
{Thương của 25 : 5 lớn hơn thương của 30 : 5.}
{\True Nếu lấy 50 : 5 thì được thương là 10.}
{\True Số chia trong bảng chia 5 luôn luôn bằng 5.}
\loigiai{
\begin{itemize}
    \item Ý a) Đúng: Đúng.
    \item Ý b) Sai: $25 : 5 = 5 < 30 : 5 = 6$.
    \item Ý c) Đúng: $50 : 5 = 10$. Đúng.
    \item Ý d) Đúng: Đúng theo tên gọi bảng chia.
\end{itemize}
\begin{itemize}
    \item a) \textbf{Đúng}.
    \item b) \textbf{Sai}.
    \item c) \textbf{Đúng}.
    \item d) \textbf{Đúng}.
\end{itemize}
}
\end{ex}

\begin{ex} %[2D8V44]
Một lớp học có 35 học sinh chia đều thành 5 nhóm:
\choiceTF
{\True Mỗi nhóm có 7 học sinh.}
{Mỗi nhóm có số học sinh là số chẵn.}
{\True Tổng số học sinh của 2 nhóm là 14 học sinh.}
{Nếu mỗi nhóm bớt đi 1 bạn thì cả lớp còn lại 25 bạn.}
\loigiai{
\begin{itemize}
    \item Ý a) Đúng: $35 : 5 = 7$ học sinh.
    \item Ý b) Sai: 7 là số lẻ.
    \item Ý c) Đúng: $7 \times 2 = 14$ học sinh.
    \item Ý d) Sai: Cả lớp bớt $5 \times 1 = 5$ bạn thì còn $35 - 5 = 30$ bạn (không phải 25).
\end{itemize}
\begin{itemize}
    \item a) \textbf{Đúng}.
    \item b) \textbf{Sai}.
    \item c) \textbf{Đúng}.
    \item d) \textbf{Sai}.
\end{itemize}
}
\end{ex}

\textbf{Phần 3. Trả lời ngắn}
\setcounterpageref{ex}{0}

\begin{ex} %[2D8V44]
Tính kết quả: 45 : 5.
\shortans{9}
\loigiai{
\textbf{Đáp số:} $9$\\
$45 : 5 = 9$.
}
\end{ex}

\begin{ex} %[2D8V44]
Có 30 lít dầu rót đều vào 5 can. Hỏi mỗi can chứa bao nhiêu lít dầu?
\shortans{6}
\loigiai{
\textbf{Đáp số:} $6$\\
$30 : 5 = 6\text{ L}$.
}
\end{ex}

\begin{ex} %[2D8V44]
Tính giá trị: 40 : 5 + 18.
\shortans{26}
\loigiai{
\textbf{Đáp số:} $26$\\
$40 : 5 = 8$; $8 + 18 = 26$.
}
\end{ex}

\begin{ex} %[2D8V44]
Một đoạn dây dài 25 cm cắt thành 5 đoạn bằng nhau. Mỗi đoạn dài bao nhiêu xăng-ti-mét?
\shortans{5}
\loigiai{
\textbf{Đáp số:} $5$\\
$25 : 5 = 5\text{ cm}$.
}
\end{ex}

\begin{ex} %[2D8V44]
Tìm x biết: x x 5 = 50.
\shortans{10}
\loigiai{
\textbf{Đáp số:} $10$\\
$x = 50 : 5 = 10$.
}
\end{ex}

\begin{ex} %[2D8V44]
Có 35 quả cam chia đều vào các đĩa, mỗi đĩa 5 quả. Hỏi chia được bao nhiêu đĩa cam?
\shortans{7}
\loigiai{
\textbf{Đáp số:} $7$\\
$35 : 5 = 7$ đĩa.
}
\end{ex}
